\ssr{ЗАКЛЮЧЕНИЕ}

В ходе выполнения лабораторной работы №2 были выполнены следующие задачи:
\begin{enumerate}
	\itemразработан рекурсивный и нерекурсивный алгоритмы решения задачи нахождения максимального числа последовательности;
	\itemописаны средства разработки и инструменты замера процессорного времени выполнения реализации алгоритмов;
	\itemреализованы разработанные алгоритмы;
	\itemпроизведено тестирование реализации алгоритмов;
	\itemвыполнена теоретическая оценка затрачиваемой реализацией каждого алгоритма памяти;
	\itemвыполнены замеры процессорного времени выполнения реализации алгоритмов в зависимости от варьируемого входа;
	\itemоценена трудоёмкость двух алгоритмов или их реализаций в худшем случае;
	\itemсравнены результаты замеров процессорного времени и оценки трудоёмкости;
	\itemсделаны выводы из сравнительного анализа реализации рекурсивного и нерекурсивного алгоритмов решения одной и той же задачи, заданной вариантом, по критериям ёмкостной эффективности (на материале теоретической оценки пиковой временной эффективности.
\end{enumerate}

В ходе лабораторной работы было выяснено, что итерационный алгоритм поиска максимально элемента массива выполняется быстрее в 7 раз в лучшем случае и в 4 раза в худшем случае чем рекурсивный алгоритм. К тому же итерационной алгоритм использует меньше памяти, потому что не происходит множественного вызова рекурсии. Кроме этого, рекурсивный алгоритм не получиться использовать на массиве большого объёма по причине переполнения стека.
